\documentclass{resume}

\usepackage{xeCJK}
\setCJKmainfont{SourceHanSerifSC-Medium}
\setCJKsansfont{SourceHanSerifSC-Medium}
\setCJKmonofont{SourceHanSerifSC-Medium}

\begin{document}

\name{林知远}
\basicInfo{
      \email{zhiyuan.lin@example.com} \textperiodcentered\
      \phone{+(86) 138-0000-0000} \textperiodcentered\
      \wechat{zhiyuanlin}\textperiodcentered\
      \homepage[zhiyuanlin.dev]{https://zhiyuanlin.dev/}
}

% ==================== 核心定位 ====================
\section{核心定位}
\textbf{云原生基础架构工程师｜5 年分布式系统 \& K8s 平台研发经验}
\begin{itemize}[parsep=0.25ex]
      \item 先后在腾讯云、蚂蚁集团、字节跳动从事基础架构研发，覆盖存储系统、容灾多活、K8s 调度、Service Mesh 全链路
      \item 主导字节 KubeFlow 智能调度平台，管理 \textbf{5000+ GPU 节点}，集群利用率从 62\% 提升至 85\%；设计蚂蚁多活容灾方案，实现机房级故障 \textbf{30 秒}自动切换
      \item 浙江大学计算机硕士（分布式系统方向），ACM-ICPC 亚洲区域赛银牌
\end{itemize}

% ==================== 教育经历 ====================
\section{教育经历}
\datedsubsection{\textbf{浙江大学 · 计算机科学与技术学院}, 硕士}{2019/09 -- 2022/06}
\begin{itemize}
      \item 计算机科学与技术，研究方向：分布式系统与云计算，发表 CCF-B 论文 1 篇
\end{itemize}
\datedsubsection{\textbf{浙江大学 · 软件学院}, 本科}{2015/09 -- 2019/06}
\begin{itemize}
      \item 软件工程，GPA 3.8/4.0
\end{itemize}

% ==================== 工作经历 ====================
\section{工作经历}

% ---------- 字节跳动 ----------
\datedsubsection{\textbf{\href{https://www.bytedance.com/}{字节跳动}}，北京}{2024/03 -- 至今}
\role{基础架构部 · 云原生团队}{高级后端开发工程师}

\textbf{项目一：KubeFlow 智能调度平台}（2024.06 - 至今 | 技术负责人）
\begin{itemize}[parsep=0.25ex]
      \item \textbf{【调度引擎】}主导调度引擎架构设计，基于 K8s Scheduler Framework 扩展自定义插件；实现 Gang Scheduling、Topology-Aware Scheduling 等策略，GPU 集群利用率从 62\% 提升至 \textbf{85\%}
      \item \textbf{【多租户管理】}设计 Quota 管理系统，支持层级配额、弹性借用、抢占优先级；构建资源画像系统，实时追踪 2000+ 团队使用模式
      \item \textbf{【智能运维】}基于 Claude Code + TDD 构建 K8s Operator，自动化处理节点故障迁移；开发智能告警聚合系统，告警噪音降低 \textbf{70\%}
      \item 技术栈: Go, Kubernetes, Scheduler Framework, Prometheus, gRPC, etcd
\end{itemize}

\textbf{项目二：服务网格控制面}（2024.03 - 2024.06 | 核心开发）
\begin{itemize}[parsep=0.25ex]
      \item \textbf{【性能优化】}实现基于 xDS 协议的配置分发引擎，支撑 \textbf{100 万+ Pod}；配置推送延迟 P99 从 800ms 降至 \textbf{200ms}
      \item \textbf{【灰度发布】}设计灰度策略引擎，支持按流量比例、用户标签、地域等多维度灰度
      \item 技术栈: Go, Envoy, xDS, gRPC, Kubernetes, Istio
\end{itemize}

% ---------- 蚂蚁集团 ----------
\datedsubsection{\textbf{\href{https://www.antgroup.com/}{蚂蚁集团}}，杭州}{2022/07 -- 2024/02}
\role{技术风险部 · 容灾架构组}{后端开发工程师}

\textbf{多活容灾平台}（保障支付宝核心链路高可用，覆盖 3 数据中心、2000+ 应用）
\begin{itemize}[parsep=0.25ex]
      \item \textbf{【流量调度】}设计基于 DNS + SDK 的多层流量调度方案，实现机房级故障 \textbf{30 秒}自动切换（RPO=0）；构建流量染色体系，支持单元化路由和灰度验证
      \item \textbf{【数据一致性】}主导跨机房数据同步方案，基于 binlog 实现准实时复制；跨城同步 P99 从 3s 降至 \textbf{500ms}，一致性达 99.999\%
      \item \textbf{【演练平台】}从零构建自动化容灾演练平台，支持串行/并行/条件分支编排；年度 200+ 次演练，发现并修复 35 个潜在缺陷
      \item 技术栈: Java (SpringBoot), OceanBase, ZooKeeper, RocketMQ, SOFA Stack
\end{itemize}

% ---------- 腾讯云 ----------
\datedsubsection{\textbf{\href{https://cloud.tencent.com/}{腾讯云}}，深圳}{2021/07 -- 2022/06}
\role{云产品部 · 对象存储组}{后端开发实习生}

\textbf{COS 对象存储服务}
\begin{itemize}[parsep=0.25ex]
      \item \textbf{【冷热分层】}参与冷热分层功能开发，设计生命周期规则引擎，冷数据存储成本降低 \textbf{40\%}
      \item \textbf{【性能优化】}优化大文件分片上传，吞吐量提升 35\%；修复 Compaction 长尾延迟，P99 读延迟降低 \textbf{60\%}
      \item 技术栈: Go, C++, LevelDB, RocksDB, gRPC, Prometheus
\end{itemize}

% ---------- MSRA ----------
\datedsubsection{\textbf{微软亚洲研究院 (MSRA)}，北京}{2020/07 -- 2020/12}
\begin{itemize}[parsep=0.25ex]
      \item 参与分布式机器学习框架优化，改进 AllReduce 通信算法，实验结果被导师采纳进后续论文
\end{itemize}

% ==================== 个人项目 ====================
\section{个人项目}

\textbf{KubeWatch — K8s 集群巡检 Operator}（\href{https://github.com/zhiyuanlin/kubewatch}{github.com/zhiyuanlin/kubewatch}）
\begin{itemize}[parsep=0.25ex]
      \item 基于 Kubebuilder 构建 Kubernetes Operator，设计 CRD 声明式巡检规则，支持资源水位、证书过期、镜像漏洞等 12 类检查
      \item 实现 Webhook 告警通知（飞书/钉钉/Slack），异常自动生成修复建议
\end{itemize}

% ==================== 技能 ====================
\section{技能}
\begin{itemize}[parsep=0.25ex]
      \item \textbf{云原生}: Kubernetes, Scheduler Framework, Kubebuilder, CRD, Helm Charts, Service Mesh (Istio/Envoy), xDS, Docker
      \item \textbf{分布式系统}: 高可用架构、容灾多活、流量调度、数据一致性（Raft/Paxos）；全链路压测、混沌工程
      \item \textbf{后端开发}: Go (主力)、Java (SpringBoot)、Python；gRPC、etcd、RocketMQ；MySQL、OceanBase、Redis
      \item \textbf{存储系统}: 对象存储（冷热分层）、LSM-Tree 引擎（LevelDB/RocksDB）、Prometheus 监控体系
\end{itemize}

% ==================== 荣誉 ====================
\section{荣誉}
\begin{itemize}
      \item \textbf{字节跳动 2024 H2 Outstanding Contributor} - 基础架构部（2024年12月）
      \item \textbf{蚂蚁集团 2023 年度技术突破奖} - 多活容灾平台项目（2023年12月）
      \item \textbf{ACM-ICPC 亚洲区域赛银牌} - 杭州站（2018年）
\end{itemize}

% ==================== 致谢 ====================
\section{致谢}
\begin{itemize}
      \item 感谢您花时间阅读我的简历，期待能有机会和您共事
\end{itemize}

\end{document}
